\documentclass{article}

\usepackage[dblblindworkshop]{neurips_2026}
\workshoptitle{SLM-Agents: 1st Workshop on SLMs for Agentic Systems}

\usepackage[utf8]{inputenc}
\usepackage[T1]{fontenc}
\usepackage{hyperref}
\usepackage{url}
\usepackage{booktabs}
\usepackage{amsmath}
\usepackage{amsfonts}
\usepackage{microtype}
\usepackage{xcolor}
\hypersetup{hidelinks}

\title{From Tokens per Second to Time to First Action:\\
Co-Designing Compact WebGPU Agents}

\author{Anonymous Author(s)}

\begin{document}

\maketitle

\begin{abstract}
An agent cannot execute a partial JSON token, yet browser language models are usually described
by decoded tokens per second.  We study \emph{time to first complete action} (TTFA), measured from
prompt tokenization through independent schema validation, and pair it with exact action success
before a deadline.  We implement a concurrency-one WebGPU benchmark and a one-forward structured
policy with route, tool-selection, and grounded-argument heads.  A matched 10.5M-parameter hybrid
backbone improves held-out bits per byte over all attention in three approximately
one-token-per-parameter training seeds (mean difference \(0.07295\), hybrid favored in 3/3), while
an exact trained checkpoint sustains 116--138 median decoded tokens/s on one Apple M5.  Our main
deployment finding is that the hidden-state extraction point is part of the model interface: a
13.0M-parameter structured policy reaches 25\,ms median TTFA but produces 0/19 required tool calls
and 0/8 executable DOM successes when fixed-length padding is inserted before the assistant
marker.  The same checkpoint routes 17/20 cases and selects 17/19 tools in the natural-marker
offline audit; moving the padding after the marker restores those counts at both 128 and 512
tokens.  Native PyTorch, fp32 ONNX, fp16 ONNX, and serialized heads agree on 20/20 normalized
actions, ruling out export/precision drift as the observed cause.  The result is a negative
deployment-methodology finding: latency, capability, and feature-contract parity must be joined on
the same artifact.  The 34M screen, matched autoregressive controls, fresh capability set, and
anonymous reproduction package remain explicitly unfinished.
\end{abstract}

\section{Why complete actions, not token rate?}

Interactive agents must return a valid decision quickly enough to preserve user flow.  Token rate
is an incomplete objective: it depends on the tokenizer, excludes prompt processing and
validation, and rewards a token that opens a JSON object as much as one that completes an action.
Browser deployment further couples model architecture to ahead-of-time WebGPU kernels,
graph/operator optimization, and paged-KV-cache sequence management \citep{ruan2026webllm}.
WebLLM uses the MLC/TVM stack rather than the ONNX Runtime Web stack measured here, so it motivates
deployment constraints rather than documenting our runtime.  At 125M and 350M parameters,
MobileLLM finds advantages from deep-and-thin blocks, embedding sharing, grouped-query attention,
and blockwise sharing \citep{liu2024mobilellm}; whether those findings extrapolate to 10--34M
parameters is a hypothesis for our matched ablation, not an imported result.

We ask a narrower systems question: \emph{what model and action representation maximize the
probability of a correct, dispatch-ready browser action before a stated deadline?}  Our
contributions are:
\begin{enumerate}
  \item a tokenizer-independent TTFA definition and deadline-conditioned exact-action metric;
  \item a browser harness that retains cold and warm latency, failures, bundle hashes, runtime
        metadata, raw normalized actions, and independent validation evidence;
  \item a matched all-attention versus periodic-attention backbone experiment and a one-forward
        structured dispatch path; and
  \item a preserved negative deployment result showing that a fast forward pass can still yield
        no useful action when the training and serving feature contracts differ, even when the
        exported graph is numerically parity-correct.
\end{enumerate}

We treat the tokenizer, assistant-marker position, padding direction, decision index, pointer
span, grounding code, and validator as one \emph{feature contract}.  This contract is a model
interface, not a page-level implementation detail.

For an autoregressive (AR) action requiring \(D\) decode steps,
\begin{equation}
  \mathrm{TTFA}_{\mathrm{AR}} =
  t_{\mathrm{tok}} + t_{\mathrm{prefill}} + (D-1)t_{\mathrm{step}}
  + t_{\mathrm{parse}} + t_{\mathrm{validate}} .
\end{equation}
The structured policy emits no serialized action tokens:
\begin{equation}
  \mathrm{TTFA}_{\mathrm{struct}} =
  t_{\mathrm{tok}} + t_{\mathrm{forward}} + t_{\mathrm{ground+validate}} .
\end{equation}
For budget \(B\), the primary score is
\[
  \mathrm{Success@}B =
  \Pr[\text{exact action} \land \text{schema valid} \land \mathrm{TTFA}\le B].
\]
Our clock currently begins at tokenization, not when a page observation becomes available; we
therefore call it \emph{harness TTFA}.  Repeated timing rows are never treated as independent
capability trials.  There is no action-independent token-rate threshold: with 200\,ms through the
first token and 20\,ms postprocessing, a 64-step action completes in 535.0, 377.5, or 325.0\,ms
at 200, 400, or 600 tokens/s, respectively.

\section{A browser-sized architecture and training contract}

\paragraph{Backbones.}
The 10.5M pilot uses a 16,384-token ByteLevel BPE vocabulary, tied embeddings, width 384, four
layers, six query heads, one KV head, SwiGLU, RMSNorm, RoPE, QK normalization, and a 2,048-token
limit.  The all-attention control has 10.547M parameters.  The 10.525M hybrid replaces layers one
and three with gated depthwise short convolutions.  A staged 34M screen repeats the comparison at
width 448 and 12 layers: all attention with FFN width 1,328 (34.276M) versus the pattern
\(\mathrm{[conv,conv,attn]}\times4\) with FFN width 1,152 (34.199M).  Because FFN allocation also
changes, this is a compound parameter-matched configuration treatment, not an isolated estimate
of ``convolution instead of attention.''

We audited disclosed frontier designs rather than shrinking them mechanically.  Kimi K3 combines
69 KDA and 24 gated-MLA layers, Attention Residuals, and Stable LatentMoE in a 2.8T-total,
104B-active-parameter MoE \citep{kimiteam2026k3}; Kimi Linear uses a uniform 3:1
KDA-linear/full-MLA layer ratio and evaluates context lengths through one million tokens
\citep{kimiteam2025linear}.  Kimi K2 reports large-scale agentic-data synthesis and joint RL in
real and synthetic environments; Kimi K2.5 adds continual mixed visual/text pretraining and
learned parallel-agent orchestration \citep{kimiteam2025k2,kimiteam2026k25}.  We transfer only
explicit stage accounting and verified parallel-task construction, not its vision, MoE/MLA, or
swarm runtime.  GLM-4.5 reports expert-model iteration and RL across agentic, reasoning, and
coding tasks \citep{glmteam2025glm45}.  SOLAR duplicates
and trims compatible pretrained layers, then continues pretraining after depth up-scaling
\citep{kim2024solar}.  The released 314B Grok-1 MoE is a raw base checkpoint and therefore
architecture evidence only \citep{xai2024grok}; the separate Grok 4.5 report describes RL on
hundreds of thousands of tasks centered on multi-step software engineering and other technical
work, plus hours-long asynchronous rollouts, but discloses no architecture
\citep{xai2026grok45}.

At 10--34M parameters and 2K context, we exclude MoE, latent attention, 100K-scale vocabularies,
and bespoke long-context kernels because resident weights, vocabulary-projection cost, and kernels
absent from our audited ONNX Runtime Web deployment path would consume the measured budget; this
is an engineering hypothesis, not evidence that those components are intrinsically inferior.
We transfer only testable
principles: periodic global mixing, reduced KV-head count, deep/thin dense baselines, stable
normalized attention, and explicit post-training stages.  The single KV head is our own
MQA-style deployment choice.  Our convolution is not KDA, and our from-scratch run is not
SOLAR-style model growth.

\paragraph{Structured dispatch.}
The final normalized hidden state feeds a text-versus-tool route gate.  Tool calls use a global
dense two-tower selector over precomputed tool-description embeddings.  Learned start/end pointer
heads ground configured strings, paths, and URLs in the natural prompt span; remaining primitives
are deterministically grounded and independently checked against JSON Schema.  The WebGPU action
graph returns hidden states only, avoiding a 16K vocabulary projection.  The exact pilot contains
10,524,544 backbone parameters and 2,499,333 active structured-head parameters (13,023,877 total);
its fixed-action payload is 27,091,138 bytes excluding page code and ONNX Runtime.

\paragraph{Data and stages.}
The pilot corpus contains 120,014,016 accepted characters from revision-pinned educational web,
synthetic-text, and permissively licensed Python sources.  After filtering and a stable document
split, a tokenizer trained on training documents only produces 28,045,897 packed training tokens
and a 240-document validation set.  Exact content hashes bind raw data, split, tokenizer, configs,
checkpoints, and scorecards.  Training follows a project-owned
tokenizer--pretrain--midtrain--SFT--RL pipeline.  The end-to-end tokenizer, pretraining, and SFT
organization is informed by nanochat \citep{karpathy2025nanochat}; LocalAgent independently adds
midtraining and bounded offline agent RL, rather than attributing those stages to nanochat's
current speedrun.  All stages use the project's canonical conversation schema.  Midtraining masks
user tokens, SFT
optimizes assistant actions and structured heads, and the bounded offline RL control uses exact
tool-AST/text rewards.  The final corpus gate additionally requires revision and license
provenance, a content-bound split map shared by every stage, train-only tokenizer fitting, and
decontamination against every frozen evaluation suite.  BFCL, BrowserGym, Mind2Web, and WebLINX
now pass strict prompt-only replay and are bound by one self-hashed four-suite manifest.  The
Mind2Web v2 ranker-bound export contains 9,378 prompts with a 1,771-byte maximum; the earlier
full-DOM v1 attempt remains a historical fail-closed result at 858,832 bytes against a
524,288-byte cap.  These exclusions are decontamination evidence, not benchmark scores.  The
independently reverified paper-v5 acquisition contains 507,082 documents and 2.200B accepted
characters.  Preparation retained 504,010 documents and packed 523,358,082 training plus
5,311,528 validation tokens with a train-only 16K tokenizer.  Independent verification reproduced
the content-addressed freeze.  Its bounded, explicitly non-exhaustive audit screened 19,334
supplied normalized denylist prompts and removed 15 documents.  Hence no native benchmark or 34M
training result is claimed.  The pretraining entry point rebuilds the freeze before model
construction and records its self-hash in checkpoint lineage; no 5-TPP arm has started.

\section{Pilot results: fast, better BPB, but not an agent}

\paragraph{Matched pretraining.}
Both 10.5M arms consume exactly 10,551,291 loss tokens per seed with identical draws.  Table
\ref{tab:quality} reports CPU-fp32 document scorecards for three prospectively designated seeds.
The primary difference is attention minus hybrid bits per byte (BPB), so positive is favorable to
the hybrid.  The mean is \(0.07295\); a Student-\(t\) interval with two degrees of freedom is
\([0.06825,0.07765]\).  This model-based interval assumes approximately normal seed effects, which
cannot be assessed at \(n=3\).  Hybrid wins 3/3, but the exact sign test is underpowered
(\(p=0.125\) one-sided; \(p=0.25\) two-sided).  This supports only promotion to a larger screen.

\begin{table}[t]
  \caption{Matched one-token-per-parameter quality proxy.}
  \label{tab:quality}
  \centering
  \small
  \begin{tabular}{@{}rrrr@{}}
    \toprule
    Seed & Attention BPB & Hybrid BPB & Difference \\
    \midrule
    2027 & 2.10301 & 2.03218 & +0.07084 \\
    2028 & 2.10478 & 2.03125 & +0.07353 \\
    2029 & 2.09968 & 2.02520 & +0.07449 \\
    \midrule
    Mean & 2.10249 & 2.02954 & +0.07295 \\
    \bottomrule
  \end{tabular}
\end{table}

\paragraph{WebGPU decoding.}
The exact seed-2026 pretrain-only hybrid checkpoint was exported as parity-gated fp16 prefill and
cache-bearing one-token decode graphs.  On one Apple M5 in Chrome 150 with ONNX Runtime Web
1.27.0, three page/session runs retained 720 measurements.  Median-of-run p50 wall rate was
137.71, 128.23, 123.78, and 116.48 tokens/s at 128, 512, 1,024, and 1,536 input tokens.  It cleared
100 tokens/s in every run/context, but a joint gate also required p95 time per output token at
most 10\,ms.  No context passed both conditions in every run.  ORT exposed a WebGPU device, but
neither adapter identity nor per-node provider placement/fallback; this is one-device latency
evidence, not agent quality.

\paragraph{Post-training and complete actions.}
A bounded seed-2027 hybrid then ran 64 midtraining steps, 320 SFT steps, and 32 offline-RL steps.
Agent held-out token accuracy rose from 3.71\% to 69.80\% during midtraining, while general loss
slightly worsened from 5.7064 to 5.7342.  SFT assistant-token accuracy rose from 67.29\% to
73.13\%, but teacher-forced all-assistant-token exactness was only 1/65.  RL obtained reward
variation in six of 128 groups and made 12 optimizer updates, yet every 53-row greedy held-out
metric was unchanged.  Moreover, the scheduled 25--50\% \emph{row} weight for agent data realized
only 1.60\% of input tokens and 1.04\% of loss tokens.  The final sampler must target and report
token shares, not imply them from row weights.  The runner now implements a deterministic,
checkpointed supervised-loss-token deficit scheduler and token-normalized gradient accumulation;
the paper configs select it, but no new full-budget result has been collected.

The fp16 structured graph was then evaluated on 20 frozen action cases and eight executable local
DOM tasks with exactly 512 input tokens.  Three browser runs produced 24.55--25.20\,ms median TTFA
and 34.30--34.80\,ms p95.  Capability nevertheless collapsed: every case abstained, giving 1/20
exact overall, 0/19 tool-required calls, and 0/8 DOM final states.  The apparent 100\% schema rate
only means that abstention was valid.

The stress harness had inserted real, unmasked BPE spaces \emph{before} the assistant marker and
then selected the final hidden state.  SFT heads were trained at the natural marker position.
Offline natural prompts recover 17/20 routes and 17/19 tool selections, whereas pre-marker filler
routes every case to text.  Appending the same filler after the marker restores 17/20 routes and
17/19 tool selections at both 128 and 512 tokens.  A preserved full-stack check compares native
PyTorch, fp32 ONNX, fp16 ONNX, exported heads, browser-identical grounding, and final normalization
at the corrected natural position: route, tool, arguments, and final action agree on 20/20 cases,
including 11/11 pointer spans.  Thus export and reduced precision are not the observed cause.  The
corrected fixed-compute runner appends filler after the marker, dispatches from the natural index,
and bounds pointer scans to the natural span, but its browser reruns await an external pre-run
timestamp.  Because these cases informed diagnosis, those reruns will measure deployment parity,
not untouched capability.

\section{Implications, limitations, and next experiment}

The pilot establishes three separable facts: a periodic-attention hybrid improves a tiny
language-model proxy; a trained 10.5M checkpoint can exceed 100 median tokens/s on one browser
device; and a 25\,ms structured forward can still fail every required tool action.  Joining the
first two while assuming the third would have produced a false ``real-time agent'' claim.  The
stronger systems lesson is that feature position, padding semantics, tokenizer bytes, graph
outputs, and grounding rules are one deployment contract.  A checkpoint without that contract is
not a reproducible agent artifact.

This remains work in progress.  The quality experiment is approximately one token per parameter,
the browser evidence uses one M5/Chrome configuration, and the current suites contain only 20 and
eight unique tasks.  The model is text/DOM-grounded rather than vision-first.  The action-trained
AR controls, fresh external/template-disjoint tasks, BrowserGym episodes, WASM arm, cross-device
runs, 34M five-token-per-parameter screen, and promoted 20-token-per-parameter comparison are
unfinished.  The next confirmatory comparison requires at least 200 fresh task IDs, paired
task-cluster bootstrap intervals, a structured-minus-AR exact-action lower bound above \(-0.02\),
and at least 15\% lower median-of-run p95 TTFA, all externally timestamped before outcomes are
inspected.  Until then, our defensible conclusion is methodological: optimize successful complete
actions under a deadline, and reject fast policies that cannot act.

Local inference can reduce server-side disclosure and make narrow assistance available without a
network connection.  The same low latency can accelerate erroneous, privacy-invasive, or
destructive browser actions.  Schema validation, abstention, sandboxed local tasks, and explicit
confirmation boundaries reduce but do not remove these risks; the present prototype is not
appropriate for unsupervised open-web deployment.

\bibliographystyle{plainnat}
\bibliography{references}

\input{checklist}

\end{document}
